\begin{textsample}{POS Dim 3 – 2011\_c6 – Score 1.00 – 2011\_c6/enade\_2011\_fis\_063ad74c0e.txt}  \label{ex:f3_pos_209}
CAPÍTULO 3 ANÁLISE TÉCNICA DA PROVA Este capítulo tem por objetivo apresentar o desempenho dos estudantes concluintes de Física ( Bacharelado e Licenciatura ) no ENADE/2011.
Para isso, foram calculadas as estatísticas básicas da prova em seu todo, bem como as estatísticas dos componentes relacionadas à Formação Geral, ao de Conhecimento Específico da Área e das questões discursivas isoladamente.
Além disso, as análises das questões objetivas do componente de Conhecimento Específico são apresentadas separadamente para concluintes do Bacharelado e da Licenciatura, já que as questões propostas para cada grupo foram diferentes.
Nas tabelas, são apresentados o tamanho da população inscrita e de presentes, e as seguintes estatísticas das notas8 : média do desempenho na prova, erro padrão da média, desvio padrão, nota mínima, mediana e nota máxima.
As estatísticas apresentadas neste capítulo contemplam o total de estudantes concluintes da área de Física ( Bacharelado e Licenciatura ) em 2011 do Brasil e, separadamente, por Grande Região.
Foram calculadas tendo -se em vista as seguintes agregações : ( a ) as Grandes Regiões e o país como um todo ; ( b ) a Categoria Administrativa ; e ( c ) a Organização Acadêmica.
Em relação aos gráficos de distribuição de notas, o intervalo considerado foi de 10 unidades, aberto à esquerda e fechado à direita, com exceção do primeiro intervalo, [ 0 ; 10 ], fechado em ambos os extremos.
Para os gráficos de distribuição das notas das questões discursivas, foram consideradas mais duas categorias : questão em branco e nota zero.
Todos os gráficos de distribuição de notas permitem a comparação dos resultados por habilitação : Bacharelado e Licenciatura. 3.1 ESTATÍSTICAS BÁSICAS DA PROVA 3.1.1 Estatísticas Básicas Gerais A Tabela 3.1 apresenta as estatísticas básicas da prova por grande Região dos estudantes concluintes de Física ( Bacharelado e Licenciatura ).
A população total de inscritos foi de 3.595.
Destes, 2.823 estiveram presentes, sendo 21,5 \% o índice de não comparecimento.
A região de maior abstenção foi a Sudeste ( 24,6 \% ), e a de menor abstenção foi a Sul ( 12,9 \% ). 8 Essas estatísticas e outras estão definidas no Capítulo 1.
A média das notas da prova como um todo ( nas seções seguintes serão analisados os componentes de Formação Geral e de Conhecimento Específico ) foi 34,8, sendo que os alunos da região Norte obtiveram a média mais baixa ( 33,0 ), e os da região Sul obtiveram a média mais alta ( 38,5 ).
As demais médias foram : 33,6 na região Nordeste ; 34,3 na região Sudeste ; e 36,4 na região Centro-Oeste.
O desvio padrão para o Brasil como um todo foi 15,1, sendo o maior desvio padrão encontrado na região Sudeste ( 17,5 ), e o menor, na região Norte ( 11,8 ), indicando uma dispersão um pouco menor das notas desta última região.
A região que obteve a maior nota máxima foi a Sul ( 81,6 ), ao passo que a região que atingiu a menor nota máxima foi a Norte ( 70,8 ).
A mediana do Brasil como um todo foi 34,7, sendo a maior mediana obtida na região Sul ( 38,9 ), e a menor obtida na Norte ( 32,5 ).
A nota mínima foi zero em todas as regiões, sem exceção.
O comportamento das notas dos estudantes de todo o Brasil pode ser observado no Gráfico 3.1 que apresenta a distribuição das mesmas em colunas diferentes para alunos de Bacharelado e de Licenciatura.
A distribuição de notas dos concluintes de Bacharelado tem a moda distribuída em dois intervalos contíguos, sendo ( 30 ; 40 ] e ( 40 ; 50 ] os intervalos modais.
Existe também um modesto máximo local no primeiro grupo de notas, [ 0 ; 10 ].
Para as notas dos estudantes de Licenciatura a distribuição é unimodal, com moda no intervalo ( 30 ; 40 ].
Os coeficientes de assimetria das distribuições das notas das duas habilitações são negativos : – 0,18 para o Bacharelado e – 0,08 para a Licenciatura.
Como tais coeficientes são pequenos observa -se que as distribuições são aproximadamente simétricas.
As distribuições por Grande Região também apresentam leve assimetria negativa.
Apenas para concluintes de Licenciatura das regiões Norte e Centro-Oeste os coeficientes de assimetria são positivos, mas também próximos de zero, 0,02 e 0,14, respectivamente.
Os Gráficos 3.2, 3.3 e 3.4 apresentam informações referentes à média da nota final dos Participantes, desagregando os resultados de acordo com, respectivamente, as Grandes Regiões do país, a Categoria Administrativa e a Organização Acadêmica.
Os gráficos apresentam o valor da média das notas como uma barra e os extremos do intervalo de confiança de 95 \% como linhas verticais unidas por uma linha horizontal na forma da letra H maiúscula.
Considerando -se o gráfico de notas segundo Grande Região ( Gráfico 3.2 ), observa -se que existe diferença estatisticamente significativa ao nível de 95 \% entre a maior média, obtida na região Sul ( 38,5 ) e as médias das regiões Sudeste ( 34,3 ), Nordeste ( 33,6 ) e Norte ( 33,0 ).
Levando -se em conta os estudantes por Categorias Administrativas ( Gráfico 3.3 ), observa -se que não existe diferença entre as médias das notas das IES Públicas e Privadas, ambas iguais a 34,8.
Tendo como foco o Gráfico 3.4, que apresenta as notas médias das provas segundo Organização Acadêmica, constata -se que existe diferença estatisticamente significativa ao nível de 95 \% entre a média das notas dos estudantes provenientes de Universidades ( 34,9 ) em relação aos de Faculdades ( 31,4 ).
Já o intervalo de confiança da média dos alunos de Centros Acadêmicos ( 32,9 ) possui interseção com os intervalos de confiança das médias dos alunos dos outros dois tipos de Organização Acadêmica.
Os intervalos de confiança para as médias dos alunos de Faculdades e Centros Universitários são maiores pelo fato de haver menos cursos de Física nestes tipos de Organizações Acadêmicas e, por tanto, menos alunos que fizeram a prova. 3.1.2 Estatísticas Básicas no Componente de Formação Geral A Tabela 3.2 apresenta as estatísticas básicas em relação ao componente da prova que avalia a Formação Geral dos estudantes concluintes.
Os alunos de todo Brasil obtiveram desempenho médio de 47,4.
A maior média foi obtida na região Centro-Oeste ( 52,3 ), e a menor, na região Sudeste ( 45,7 ).
As demais médias foram : 49,8 na região Norte ; 46,4 na região Nordeste ; e 48,0 na região Sul.
Quanto à variabilidade, o desvio padrão das notas dos estudantes do Brasil como um todo foi 21,4.
O maior desvio padrão foi obtido na região Sudeste ( 23,6 ), e o menor, na região Norte ( 18,7 ).
Os demais desvios padrões foram : 20,4 na região Nordeste ; 20,7 na região Sul ; e 19,4 na região Centro-Oeste.
A maior nota no componente de Formação Geral da prova do ENADE foi 100,0, obtida por pelo menos um aluno da região Norte, enquanto a menor nota máxima foi obtida na região Sul ( 90,5 ).
Nas outras regiões as notas máximas foram : 91,5 na região Nordeste, 94,0 na região Sudeste ; e 99,0 na região Centro-Oeste.
A mediana do Brasil como um todo foi 50,0, sendo a menor mediana encontrada na região Nordeste ( 49,0 ), e a maior encontrada na região Centro-Oeste ( 55,5 ).
A nota mínima nesta parte foi zero em todas as regiões, sem exceção.
O Gráfico 3.5 propicia a avaliação do desempenho dos estudantes no componente de Formação Geral a partir do histograma da distribuição das notas correspondentes ao Bacharelado e à Licenciatura.
Na distribuição de notas dos alunos de Licenciatura a moda é dividida em dois intervalos contíguos, ( 40 ; 50 ] e ( 50 ; 60 ].
Existe também um modesto máximo local no primeiro grupo de notas, [ 0 ; 10 ].
A distribuição das notas dos estudantes do Bacharelado tem pontos de máximo nos intervalos [ 0 ; 10 ] e ( 50 ; 60 ].
O máximo local no primeiro grupo compreende mais de 15 \% dos alunos de bacharelado.
Nota -se, ainda, que no Gráfico 3.5 as notas apresentam uma maior dispersão do que no Gráfico 3.1 ( distribuição das notas da prova ), confirmado pela comparação dos desvios padrões : 15,1 para a nota da prova como um todo e 21,4 para o componente de Formação Geral.
Para o componente de Formação Geral, o coeficiente de assimetria da distribuição das notas dos estudantes, como na prova como um todo, também é negativo ( – 0,52 ).
Desagregando por habilitação observa -se que o coeficiente da distribuição dos concluintes da Licenciatura ( – 0,53 ) tem maior valor absoluto do que o do Bacharelado ( – 0,33 ), o que indica uma assimetria à esquerda mais acentuada para o primeiro grupo.
A menos dos alunos de Bacharelado em Física da região Norte ( coeficiente 0,11 ), tomando -se os alunos em conjunto ou por habilitação, em todas as Grandes Regiões os histogramas possuem assimetria negativa.
Nos Gráficos 3.6, 3.7 e 3.8 são apresentadas as informações referentes ao desempenho dos Concluintes no componente de Formação Geral, em diferentes agregações : Grande Região do país, Categoria Administrativa e Organização Acadêmica.
Observa -se pelo Gráfico 3.6 que existe diferença estatisticamente significativa entre a maior média das notas no Componente de Formação Geral, obtida na região Centro-Oeste ( 52,3 ), e as médias das regiões Sul ( 48,0 ), Sudeste ( 45,7 ) e Nordeste ( 46,4 ).
A diferença só não é significativa em relação à segunda maior média, a da região Norte ( 49,8 ).
No Gráfico 3.7, que representa as notas médias no Componente de Formação Geral segundo Categoria Administrativa do país, observa -se que existe diferença estatisticamente significativa entre as médias, diferentemente do que ocorreu para as notas da prova como um todo.
Para o componente de Formação Geral, os concluintes de Física ( Bacharelado e Licenciatura considerados em conjunto ) das IES Públicas ( 46,4 ) obtiveram uma média menor do que os das IES Privadas ( 53,3 ).
Considerando -se o tipo de Organização Acadêmica, nota -se, no Gráfico 3.8, que não existe diferença estatisticamente significativa entre as médias dos alunos de Universidades ( 47,2 ), Centros Universitários ( 52,6 ) e Faculdades ( 50,4 ). 3.1.3 Estatísticas Básicas do Componente de Conhecimento Específico A Tabela 3.3 apresenta as estatísticas básicas referentes ao componente de Conhecimento Específico da área de Física ( Bacharelado e Licenciatura ).
A média do desempenho dos alunos do Brasil como um todo foi 30,5.
A maior média foi obtida na região Sul ( 35,4 ), e a menor, na região Norte ( 27,4 ).
As demais médias foram : 29,3 na região Nordeste ; 30,5 na região Sudeste ; e 31,1 na região Centro-Oeste.
Quanto à variabilidade das notas, o desvio padrão do Brasil como um todo foi 15,9, sendo o maior desvio padrão observado na região Sudeste ( 18,1 ) e o menor na região Norte ( 12,9 ).
Os demais desvios foram : 14,5 na região Nordeste ; 16,6 na região Sul ; e 13,0 na região Centro-Oeste.
A mediana das notas dos estudantes de todo o Brasil foi 29,6.
A maior mediana ocorreu na região Sul ( 33,2 ) e a menor na região Norte ( 25,8 ).
As demais medianas foram : 27,6 na região Nordeste ; 29,6 nas regiões Sudeste e Centro-Oeste.
A nota máxima do Brasil como um todo foi 86,9, sendo obtida por pelo menos um aluno da região Sudeste.
As demais notas máximas foram : 68,1 na região Norte ; 76,6 na região Nordeste ; 86,9 na região Sul ; e 74,7 na região Centro-Oeste.
A nota mínima foi zero em todas as regiões.
Assim como os Gráficos 3.1 e 3.5, o Gráfico 3.9, apresentado a seguir, proporciona uma avaliação do desempenho de concluintes dos cursos de Bacharelado e de Licenciatura em relação ao componente de Conhecimento Específico com um histograma da distribuição das notas correspondentes.
Dentre as distribuições apresentadas, estas são as mais concentradas em notas baixas.
A distribuição das notas dos concluintes do Bacharelado e da Licenciatura são unimodais, com intervalo modal ( 20 ; 30 ] para o primeiro e dividido em dois intervalos contíguos, ( 20 ; 30 ] e ( 30 ; 40 ], para a segunda habilitação.
O coeficiente de assimetria da distribuição das notas do componente de Conhecimento Específico é positivo tanto para o curso de Bacharelado ( 0,15 ) quanto para Licenciatura ( 0,37 ).
Nota -se pelo histograma ( Gráfico 3.9 ) que para ambas as habilitações há cauda mais longa à direita.
Considerando -se as distribuições de notas nas diferentes regiões, apenas para concluintes de Bacharelado das regiões Norte e Centro-Oeste o coeficiente de assimetria é negativo, respectivamente, – 0,18 e – 0,11.
Para todas as outras agregações por regiões e habilitação os coeficientes de assimetria são positivos.
A assimetria positiva mais acentuada é a da distribuição de notas dos alunos de Licenciatura da região Centro-Oeste ( 0,61 ) e a mais leve é dos alunos de licenciatura da região Sul ( 0,17 ).
Os Gráficos 3.10, 3.11 e 3.12 apresentam uma comparação dos resultados em relação à Grande Região do país, à Categoria Administrativa e à Organização Acadêmica, agora levando em conta o desempenho de alunos no componente de Conhecimento Específico.
Pelo Gráfico 3.10, observa -se que existe diferença estatisticamente significativa entre a maior média das notas no componente de Conhecimento Específico, da região Sul ( 35,4 ) e as médias das demais regiões.
Existe também diferenças estatisticamente significativas entre a média da região Norte ( 27,4 ) e das regiões Centro-Oeste ( 31,1 ) e Sudeste ( 30,5 ).
Quanto à Categoria Administrativa ( Gráfico 3.11 ), observa -se que existe diferença estatisticamente significativa entre as médias, e a média das instituições Públicas ( 30,9 ) é maior do que das Privadas ( 28,6 ).
Este comportamento é diferente do encontrado para a parte de Formação Geral ( Gráfico 3.7 ), no qual a média das IES Privadas era significativamente maior.
Também difere do comportamento encontrado para a prova como um todo ( Gráfico 3.3 ) que não apresentou diferença estatisticamente significativa entre as médias dos dois tipos de Categorias Acadêmicas.
Quanto ao Gráfico 3.12, observa -se que existe diferença estatisticamente significativa ao nível de 95 \% entre as notas do componente de Conhecimento Específico dos estudantes oriundos das Universidades ( 30,8 ) em relação aos de Centros Universitários ( 26,3 ) e Faculdades ( 25,0 ).
Tal resultado difere do encontrado para a parte de Formação Geral da prova, onde não havia diferença estatisticamente significativa. 3.2 ANÁLISE DAS QUESTÕES OBJETIVAS 3.2.1 Componente de Formação Geral A Tabela 3.4 apresenta as estatísticas básicas relativas às oito questões objetivas do componente da prova que abrange a Formação Geral dos estudantes.
A média do Brasil foi 50,2.
A menor média foi encontrada na região Sudeste ( 47,7 ) e a maior na região Centro-Oeste ( 52,9 ).
As demais médias foram : 52,8 na região Norte ; 50,1 na região Nordeste ; e 52,0 na região Sul.
O desvio padrão do Brasil foi 22,7, sendo o maior desvio padrão encontrado na região Sudeste ( 24,6 ) e o menor na região Centro-Oeste ( 20,4 ).
Os demais desvios foram : 21,7 nas regiões Norte e Nordeste e 22,1 na região Sul.
As medianas ( 50,0 ), as notas máximas ( 100,0 ) e as notas mínimas ( 0,0 ) foram iguais para todas as regiões.
A Tabela 3.5 apresenta os índices de facilidade e discriminação ( ponto bisserial ) das questões objetivas do componente de Formação Geral.
Quanto ao índice de facilidade, foram usadas as seguintes cores para diferenciar o nível de dificuldade da questão : • Azul para as questões classificadas com índice muito fácil ( índice \&gt ; =0,86 ), verde para as questões classificadas com índice fácil ( 0,61 a 0,85 ), amarelo para as questões classificadas com médio ( 0,41 a 0,60 ), vermelho para as questões classificadas com difícil ( 0,16 a 0,40 ) e roxo para as questões classificadas com muito difícil ( \&lt ; =0,15 ).
Já quanto ao índice de discriminação, foram usadas as seguintes cores para qualificar a questão : • As questões classificadas com índice fraco receberam a cor vermelho ( índice \&lt ; =0,19 ), as classificadas com médio receberam a cor amarelo ( 0,20 a 0,29 ), as classificadas com bom receberam a cor verde ( 0,30 a 0,39 ) e as classificadas com muito bom ( \&gt ; =0,40 ) receberam a cor azul.
As questões objetivas do componente de Formação Geral, segundo o índice de facilidade, foram assim avaliadas : das oito questões, nenhuma teve o índice de facilidade classificado como muito fácil.
Quatro questões foram tidas como fáceis, por terem índice de acertos situado na faixa entre 0,61 e 0,85 ( de 61,0 \% a 85,0 \% de acertos ).
Duas questões foram consideradas de dificuldade média, situando -se no intervalo entre 0,41 e 0,60 do índice de facilidade, ou seja, houve entre 41,0 \% e 60,0 \% de acertos.
Outras duas questões foram classificadas na categoria difícil, situando -se no intervalo entre 0,16 e 0,40.
Por fim, nenhuma das questões apresentou menos de 15 \% de acertos, razão pela qual seriam classificadas como muito difíceis.
Como já comentado, para análise das questões objetivas relativas à Formação Geral segundo o poder de discriminação, utilizou -se, o índice de discriminação ponto bisserial.
Nesta análise sete das oito questões apresentaram índices acima ou igual a 0,40 e, assim, foram classificadas com índice muito bom para esse grupo de estudantes.
Uma questão teve índice de discriminação médio, entre 0,20 e 0,29, para esse grupo de estudantes.
Nenhuma questão teve nível fraco de discriminação.
O índice de facilidade variou de 0,18 a 0,70, e o de discriminação, de 0,28 a 0,64.
As sete questões com índices de discriminação muito bom, as de números 1, 2, 3, 4, 5, 6 e 7, figuraram entre os diversos níveis de dificuldade desse conjunto : quatro classificadas na categoria fácil ( questões 1, 3, 5 e 6 ) do índice de facilidade, duas na categoria médio ( questões 2 e 4 ) e uma na categoria difícil ( questão 7 ).
Em particular, a questão 3 foi a que apresentou o maior poder discriminatório, com índice 0,64, e foi também a mais fácil, com uma proporção de 0,70 acertos.
As questões 2 e 4 foram acertadas por 44 \% e 41 \% dos alunos, respectivamente, e obtiveram índice muito bom de discriminação ( respectivamente 0,50 e 0,41 ).
Já a questão 8 obteve índice de discriminação médio, 0,28, e seu índice de facilidade foi difícil ( 0,18 ).
O Gráfico 3.13, para exemplificar, analisa o comportamento da questão de número 3 de Formação Geral.
Trata -se da questão mais fácil e a que obteve o maior índice de discriminação dessa parte da prova.
Neste gráfico, cada uma das cinco curvas representa o percentual de respostas em determinada alternativa da questão, em função do número de acertos dos estudantes nesta parte da prova ( Formação Geral/Múltipla Escolha ), antes de possíveis eliminações pelo critério do ponto bisserial.
Em princípio, a soma das escolhas possíveis deveria ser igual a 100 \%, no entanto, esta soma não é alcançada o percentual dos alunos com zero acertos não aparece no gráfico, bem como o percentual daqueles que deixaram a questão em branco ou assinalaram múltiplas respostas.
Como foram oito as questões, os valores variam de 0 a 8 acertos.
A curva em vermelho corresponde à alternativa E, a correta para esta questão.
Assim, observa -se que entre os estudantes com menor número de acertos, nessa parte do exame, a situação mais frequente foi a escolha de uma das alternativas incorretas, a alternativa A ( em azul ).
Na medida em que o número de acertos aumenta, indicando desempenho melhor nesta parte da prova, aumenta concomitantemente a proporção de estudantes que selecionaram a alternativa correta E, atingindo 100 \% para os estudantes com 8 acertos.
Essa análise permite verificar como a questão discriminou os grupos de desempenho, justificando o alto índice obtido na questão.
Os gráficos relativos às demais questões de Formação Geral constam do Anexo I. 3.2.2 Componente de Conhecimento Específico – Física Bacharelado A Tabela 3.6 apresenta as estatísticas básicas em relação às questões objetivas do componente de Conhecimento Específico da prova de Física – Bacharelado, por Grande Região.
A média do Brasil deste componente foi de 37,6.
A menor média foi observada na região Sudeste ( 36,0 ) e a maior, na região Centro-Oeste ( 46,7 ).
O desvio padrão de todo o Brasil foi 22,8, sendo o menor desvio padrão encontrado na região Centro-Oeste ( 15,4 ), e o maior, na região Sudeste ( 25,1 ).
A mediana de todo o Brasil foi 39,1, a mesma encontrada na região Sul.
Nas regiões Nordeste e Sudeste a mediana foi menor, 34,8.
Já a região Centro-Oeste foi onde ocorreu a maior mediana, 47,8.
A nota máxima da prova foi 91,3, obtida nas questões objetivas do componente de Conhecimento Específico por pelo menos um aluno da região Sul.
A menor nota máxima foi 73,9, obtida na região Centro-Oeste.
A Tabela 3.7 apresenta os índices de facilidade e discriminação ( ponto bisserial ) das questões objetivas do componente de Conhecimento Específico para os estudantes de Física – Bacharelado.
Para facilitar a diferenciação das questões usaram -se as mesmas cores das Tabelas 3.5 para as diferentes classificações dos índices de facilidade e de discriminação.
As questões objetivas de Conhecimento Específico para os alunos de Bacharelado foram as de números 9 a 25 e 36 a 45.
Destas 27 questões, quatro foram anuladas pela Comissão : as de números 18, 22, 36 e 45.
Desse modo, a classificação quanto ao índice de facilidade foi estabelecida com base nas 23 questões válidas.
A partir dos índices obtidos, observa -se que mais da metade das questões objetivas da prova foi considerada difícil : 14 das 23 questões válidas.
Apenas uma questão teve índice de facilidade que a classifica como fácil e as oito demais questões foram classificadas como de dificuldade média, índice de facilidade entre 0,41 e 0,60.
Não houve questão classificada como muito difícil nem como muito fácil.
Já quanto aos índices de discriminação das questões objetivas do componente de Conhecimento Específico da prova, tem -se como resultado a seguinte classificação : 19 das 23 questões válidas foram consideradas como muito boas, enquanto três delas tiveram índice de discriminação bom e apenas uma teve índice médio.
Assim, para a maioria das questões – 21 em 23 – os índices de discriminação foram bons ou muito bons.
Constata -se, assim, que a prova – no que se refere ao componente de Conhecimento Específico – possuía boa capacidade de discriminar entre aqueles que dominam ou não o conteúdo.
Dentre as 19 questões que alcançaram índice de discriminação muito bom, situando -se no intervalo de 0,40 a 0,50 do índice, a maioria delas ( questões 9, 10, 11, 19, 20, 21, 37, 38, 40, 43 ) foi classificada na categoria difícil, quanto ao índice de facilidade.
A questão 14 foi considerada fácil e outras oito questões tiveram índice de facilidade médio ( questões 15, 16, 17, 23, 24, 25, 39 e 41 ).
A questão de número 12 foi a mais difícil dentre as 23 questões específicas válidas, com baixo índice de facilidade, apenas 19 \% de acertos.
Essa questão foi também a que apresentou poder discriminatório mais baixo, 0,25, o que comprova ter sido esta a mais difícil para os estudantes.
Destaca -se, também, a questão 21, com índice de facilidade 0,20, o que, em termos percentuais, corresponde a 20,0 \% de estudantes que responderam acertadamente.
No entanto, o índice de discriminação desta questão foi muito bom 0,51.
Como nenhuma questão teve índice de discriminação fraco, não houve questão eliminada do cômputo da nota final pelo critério ponto bisserial.
A título de exemplo das análises do comportamento das questões objetivas, o Gráfico 3.14 analisa a questão 41 do componente de Conhecimento Específico do Bacharelado.
Esta foi que obteve o maior índice de discriminação, 0,65.
Seu índice de facilidade foi médio, igual a 0,45, ou seja, 45 \% dos estudantes assinalaram a opção A, gabarito da questão.
Neste gráfico, cada uma das cinco curvas representa o percentual de respostas em determinada alternativa da questão 41, em função do número de acerto dos estudantes nesta parte da prova, antes de possíveis eliminações de questões pelo critério do ponto bisserial.
A alternativa correta A, representada no gráfico pela curva em azul, foi escolhida em maiores proporções pelos alunos com desempenho melhor nesta parte da prova.
Já as alternativas incorretas, também denominadas distratores, foram selecionadas principalmente por aqueles com notas mais baixas.
Neste caso também a soma não é sempre 100 \% por causa das questões não respondidas ou com mais de uma opção marcada.
Aqueles com nota zero, na sua quase totalidade deixaram esta questão em branco ou marcaram mais de uma alternativa, comportamento considerado inválido.
A proporção de alunos que selecionou a resposta correta A aumenta gradativamente, chegando a atingir 100 \% para mais de 18 acertos, enquanto a proporção dos que escolheram alternativas incorretas decai, a partir de um ou dois acertos, como função do número de acertos nesta parte da prova.
Os gráficos relativos às demais questões do componente de Conhecimento Específico constam do Anexo I. 3.2.3 Componente de Conhecimento Específico – Física Licenciatura A Tabela 3.8 apresenta as estatísticas básicas em relação às questões objetivas do componente de Conhecimento Específico da prova de Física - Licenciatura, por Grande Região.
A média do Brasil deste componente foi de 34,4.
A menor média foi observada na região Norte ( 31,5 ) e a maior na região Sul ( 40,8 ).
O desvio padrão de todo o Brasil foi 16,4, sendo o maior desvio padrão encontrado na região Sudeste ( 18,1 ), e o menor, na região Norte ( 14,2 ).
A mediana de todo o Brasil foi 34,8, também encontrada para as regiões Sudeste e Centro-Oeste.
A maior mediana ocorreu na região Sul ( 39,1 ) e a menor nas regiões Norte e Nordeste ( 30,4 ).
A nota máxima da prova foi 91,3, obtida por pelo menos um concluinte de Licenciatura da região Sul, nas questões objetivas do componente de Conhecimento Específico.
A menor nota máxima foi 73,9, obtida na região Norte.
A Tabela 3.9 apresenta os índices de facilidade e discriminação ( ponto bisserial ) das questões objetivas do componente de Conhecimento Específico da Licenciatura em Física.
Para facilitar a diferenciação das questões usaram -se as mesmas cores das Tabelas 3.5 e 3.7 para as diferentes classificações dos índices de facilidade e de discriminação.
As questões objetivas de Conhecimento Específico para os alunos de Licenciatura foram as de números 9 a 35.
Dentre as 27 questões objetivas da parte da prova relativa ao componente de Conhecimento Específico, duas foram anuladas pela Comissão ( questões 18 e 22 ).
Desse modo, a classificação quanto ao índice de facilidade foi estabelecida com base em 25 questões.
A partir dos índices obtidos, pode -se concluir que mais da metade das questões objetivas da prova foi considerada pelo menos difícil : das 25 questões, 16 foram classificadas como difíceis e uma como muito difícil.
Não houve questão classificada como muito fácil, ao passo que uma foi classificada como fácil, na faixa de 0,61 a 0,85 do índice de facilidade, e as outras sete consideradas médias, entre 0,41 e 0,60.
Já quanto aos índices de discriminação das questões objetivas do componente de Conhecimento Específico da prova, tem -se como resultado a seguinte classificação : 13 das 25 questões foram consideradas como boas, enquanto seis delas tiveram índice de discriminação muito bom.
Assim, para a maioria das questões – 19 em 25 – os índices de discriminação foram bons ou muito bons.
Dentre as demais, quatro foram classificadas como médias e outras duas como fracas, sendo seis, por conseguinte, a quantidade de questões nos dois patamares mais baixos de discriminação.
Constata -se, assim, que a prova – no que se refere ao componente de Conhecimento Específico – possuía capacidade boa de discriminar entre aqueles que dominam ou não o conteúdo.
Dentre as questões que alcançaram os maiores índices de discriminação, as de números 14, 15, 24, 25, 27 e 31 foram classificadas com índice muito bom, situando -se no intervalo de 0,42 a 0,50 do índice.
Dentre estas, uma ( questão 27 ) foi classificada na categoria fácil, três ( questões 14, 15 e 31 ) como médias e outras duas ( questões 24 e 25 ) como difíceis, quanto ao índice de facilidade.
A questão de número 21 foi a mais difícil dentre as 25 questões específicas válidas, com baixo índice de facilidade, apenas 13,0 \% de acertos.
Essa questão apresentou poder discriminatório médio, 0,28, o que indica que apesar de muito difícil a questão discriminou razoavelmente os alunos por desempenho.
Destacam -se, também, as questões 12 e 26 que tiveram índice de discriminação fraco.
Estas duas questões tiveram índice de facilidade baixo, ou seja, estão entre as consideradas difíceis, com índice de facilidade 20,0 \% para a questão 12 e 19,0 \% para a questão 26.
Pelo critério do índice de discriminação ponto bisserial as questões 12 e 26 foram eliminadas do cômputo da nota final.
A título de exemplo das análises do comportamento das questões objetivas, o Gráfico 3.15 analisa a questão 14 do componente de Conhecimento Específico da Licenciatura.
Esta questão foi respondida acertadamente por 56,0 \% dos estudantes apresentando índice de facilidade 0,56.
Seu índice de discriminação foi igual a 0,47, classificado como muito bom e o maior valor alcançado por este índice para o conjunto de perguntas em análise.
Neste gráfico, cada uma das cinco curvas representa o percentual de respostas em determinada alternativa da questão 14, em função do número de acertos dos estudantes nesta parte da prova, antes de possíveis eliminações de questões pelo critério do ponto bisserial.
A alternativa correta B, representada no gráfico pela curva em verde, foi escolhida em maiores proporções pelos alunos com desempenho melhor nesta parte da prova.
Já as alternativas incorretas, também denominadas distratores, foram selecionadas principalmente por aqueles com notas mais baixas.
Aqueles com nota zero, na sua quase totalidade deixaram esta questão em branco ou marcaram mais de uma alternativa, comportamento considerado inválido.
A proporção de alunos que selecionou a resposta correta B aumenta gradativamente, chegando a atingir 100 \% para os que tiveram 17 acertos ou mais, enquanto a proporção dos que escolheram alternativas incorretas decai, a partir de três acertos, como função do número de acertos nesta parte da prova.
Os gráficos relativos às demais questões do componente de Conhecimento Específico constam do Anexo I. 3.3 ANÁLISE DAS QUESTÕES DISCURSIVAS 3.3.1 Componente de Formação Geral As análises dos resultados de desempenho dos estudantes de Física ( Bacharelado e Licenciatura ) nas duas questões discursivas relativas à Formação Geral encontram -se na Tabela 3.10 e no Gráfico 3.16.
Na tabela 3.10 observa -se que a nota média nesse conjunto de questões foi inferior a obtida nas objetivas.
Os estudantes de todo o Brasil obtiveram, em Formação Geral, média 50,2 nas questões objetivas e 43,3 nas questões discursivas.
Pode -se notar, também, um aumento do desvio padrão de 22,7 nas questões objetivas do componente de Formação Geral dos alunos de todo o Brasil, para 31,0 nas questões discursivas do mesmo componente.
A maior média foi obtida na região Centro-Oeste ( 51,4 ) e a menor na região Nordeste ( 40,8 ).
A mediana de todo o Brasil, neste componente, foi 47,5.
A menor mediana foi obtida na região Nordeste ( 45,0 ) e a maior, na região Centro-Oeste ( 55,0 ).
A nota máxima ( 100,0 ) e a nota mínima ( 0,0 ) foram as mesmas em todas as regiões do Brasil.
O Gráfico 3.16 representa a distribuição das notas nas questões discursivas no componente de Formação Geral segundo a opção de Bacharelado/Licenciatura.
A moda destas distribuições ocorrem no intervalo [ 0 ; 10 ], intervalo que inclui as notas zero e a frequência de alunos que deixaram esta questão em branco.
Cumpre notar que a proporção de estudantes neste primeiro intervalo de notas é maior entre os de Bacharelado do que entre os de Licenciatura.
O coeficiente de assimetria da distribuição de notas dos concluintes do Bacharelado é positivo ( 0,24 ) enquanto o da distribuição dos concluintes de Licenciatura é negativo ( – 0,21 ).
Esse comportamento da assimetria à direita para o Bacharelado e à esquerda para a Licenciatura se repete para todas as regiões, sugerindo a repetição dos padrões dos histogramas.
A exceção ocorre para a região Centro-Oeste, onde ambos os coeficientes são negativos, ou seja, as duas distribuições possuem assimetria à esquerda.
Na sequência, os resultados verificados para cada uma das questões discursivas de Formação Geral serão apresentados, estabelecendo -se relações com os conteúdos abordados em cada uma delas.
Os comentários da Banca de docentes corretores a respeito do observado na correção das respostas dos estudantes, suas impressões e conclusões serão apresentados junto à análise de cada questão.
Cumpre esclarecer que, tendo em vista que as questões discursivas de Formação Geral são padronizadas, ou seja, constam de todas as provas, os comentários da Banca são os mesmos para todas as carreiras acadêmicas, sendo direcionados a todos os estudantes que participaram do ENADE/2011.
A seguir, serão analisados os desempenhos da Área de Física ( Bacharelado e Licenciatura ) nas duas questões discursivas de Formação Geral do ENADE/2011, comparando os resultados obtidos com comentários para cada questão. 3.3.1.1 Análise da Questão Discursiva 1 do Componente de Formação Geral Os dados de Física ( Bacharelado e Licenciatura ), obtidos a partir das respostas à questão 1, encontram -se na Tabela 3.11 e no Gráfico 3.17.
Nessa questão – de melhor desempenho dentre as duas de Formação Geral – os alunos de todo Brasil tiveram média, 49,5.
A maior média para a questão 1 foi obtida na região Centro-Oeste ( 56,9 ), e a menor, na região Nordeste ( 45,5 ).
Quanto à variabilidade das notas, o desvio padrão de todo o Brasil foi 38,3.
O menor desvio padrão foi obtido na região Centro-Oeste ( 35,5 ), e o maior desvio padrão foi obtido na região Sul ( 39,3 ).
As medianas das regiões Sudeste e do Brasil como um todo foram iguais ( 50,0 ).
Nas demais regiões as medianas foram : 55,0 nas regiões Norte e Sul ; 40,0 na região Nordeste ; e 60,0 na região Centro-Oeste.
As notas máximas e mínimas da questão discursiva 1 foram as mesmas para todas as regiões do Brasil, respectivamente, 0,0 e 100,0.
O Gráfico 3.17 mostra a distribuição das notas na questão discursiva 1 do componente de Formação Geral para alunos do Bacharelado e da Licenciatura.
Observa -se que a maior frequência para os alunos de Bacharelado corresponde aos alunos que deixaram essa questão em branco.
Para os alunos da Licenciatura a moda ocorre no intervalo ( 90 ; 100 ], apesar de também observar -se que na categoria “ em branco ” se encontra a segunda maior concentração da distribuição das notas.
O coeficiente de assimetria é negativo e muito próximo de zero para a habilitação em Licenciatura ( – 0,08 ), enquanto é positivo para o Bacharelado ( 0,28 ).
Para a Licenciatura, os coeficientes de assimetria por região são negativos, com exceção da região Nordeste.
No entanto, todos são próximos de zero.
Para o Bacharelado os coeficientes são positivos e variam de 0,05 a 0,47 por região, sendo negativo apenas para a região Centro-Oeste ( – 0,58 ). 3.3.1.2 Comentários sobre a correção das respostas à Questão Discursiva 1 De maneira geral, a aplicação da chave de correção da questão 1 de Formação Geral não apresentou qualquer dificuldade digna de menção.
Para isso contribuíram, indubitavelmente, os ajustes feitos após a correção da amostra e a reunião entre todos os membros da Banca de docentes corretores.
As poucas dúvidas, todas pontuais, apresentadas pelos corretores, foram acompanhadas e respondidas pela coordenação e subcoordenação da correção das questões de Formação Geral, por meio da ferramenta de Gerenciamento de Dúvidas do Sistema de Correção On-line.
Não houve registro de qualquer ocorrência que pusesse em xeque o padrão de resposta ou a efetividade e a adequação da chave de correção.
Explica -se : trata -se de questão com comando claro, direto e objetivo ( solicitava -se, basicamente, três vantagens justificadas de cursos a distância ), cujas respostas foram corrigidas por meio da aplicação de um chave de correção testada e aprovada previamente.
Havia absoluta clareza quanto aos critérios de avaliação da correspondência entre as respostas dos estudantes e as possibilidades de vantagens de cursos a distância admitidas como corretas no padrão de resposta oficial, além de gradações explícitas ( e fáceis de aplicar ) dos diferentes níveis de pontuação previstos.
Felizmente, portanto, não há reparo a registrar em relação à facilidade de aplicação do padrão de resposta e da chave de correção, e nem em relação à atribuição dos diferentes níveis de pontuação previstos.
Todas as dificuldades que poderiam ter obstado a correta aplicação do padrão de resposta oficial e da respectiva chave de correção foram evitadas por meio dos ajustes feitos após a correção da amostra e fartamente debatidos com toda a Banca.
Digno de nota é que quantidade tão significativa de profissionais envolvidos na mesma tarefa – tanto para a questão 1 quanto para a 2 – tenha apresentado tão poucas dificuldades na execução da correção, em termos proporcionais.
Em suma, a correção da questão 1 da prova de Formação Geral do ENADE 2011 foi exemplarmente bem planejada, servindo -se de padrão de resposta muito bem adequado à questão proposta.
Quanto ao tema desta questão, em particular, a Banca verificou que uma parcela significativa de estudantes evocou experiências bastante concretas e próximas de sua realidade.
Houve várias respostas que indicavam uma vivência pessoal de ensino superior na modalidade Educação a Distância ( EaD ), evidenciando o tom de depoentes nos textos apresentados.
Os estudantes, em sua grande maioria, utilizaram parte considerável do espaço de 15 linhas disponíveis para a resposta – e outra parcela menos significativa dos que não o fizeram demonstraram notável capacidade de atender ao comando da questão de maneira objetiva, curta e, via de regra, correta.
Registre -se, ainda, que foram relativamente poucos os casos de respostas que tenham passado ao largo do tema em pauta na questão.
As capacidades de leitura, de compreensão do comando proposto e de expressão escrita que os estudantes avaliados na edição 2011 do ENADE foram satisfatórias.
Não obstante, seria leviano perder de vista que a qualidade dos textos redigidos em resposta às questões discursivas do Exame ainda está muito aquém do que se espera de concluintes de cursos de ensino superior de todas as regiões do país.
Quanto ao conteúdo das respostas, a Banca constatou boa capacidade, por parte da maioria dos estudantes, de compreensão do tema e do comando da questão.
Foram relativamente poucos os casos de respostas que deixaram de enumerar vantagens da modalidade EaD, e proporcionalmente escassos os estudantes que citaram vantagens não previstas no padrão de resposta.
Foi frequente, a tentativa direta de atender ao comando da questão.
Os erros mais comuns, em relação ao padrão de respostas e à grade de correção, foram fruto do desdobramento em vários “ itens ” daquilo que, de acordo com o padrão de respostas oficial, representava uma única vantagem.
Destacaram -se, neste caso, as respostas que apontavam a flexibilidade de horário e/ou local como duas vantagens distintas daquela modalidade de ensino.
Quanto aos diferentes níveis de pontuação previstos, a maior causa de baixas pontuações foi a ausência de justificativas, e mesmo de argumentação, para uma ou mais das vantagens enumeradas.
Isso demonstra não apenas a objetividade das respostas, coerente com a objetividade do comando da questão ( “ enumere três vantagens de um curso a distância ” ), mas também certa dificuldade de formulação plena de um texto, ou ao menos de parágrafos, em formato dissertativo – mesmo diante de uma média de 5 linhas disponíveis para cada vantagem a enumerar.
Os acertos mais comuns, ou seja, os “ itens ” do padrão de respostas mais frequentemente mencionados foram : ( 1 ) a flexibilidade de horário e/ou local ; ( 2 ) a capilaridade do ensino a distância ; ( 3 ) a democratização do acesso à educação de qualidade ; e ( 4 ) os custos menores que os de cursos presenciais.
Dentre as vantagens previstas no padrão de respostas que foram menos citadas, destacam -se a inclusão de pessoas com comprometimento motor, a qualificação de professores e a troca de experiências entre os participantes.
Foram muito comuns, no entanto, as menções à supostamente maior facilidade de acesso a professores ou/e tutores em cursos superiores a distância. 3.3.1.3 Análise da Questão Discursiva 2 do Componente de Formação Geral A Tabela 3.12 mostra que o desempenho dos estudantes na questão 2 ( média 37,0 ) foi inferior ao obtido na questão de número 1 ( média 49,5 ).
A região Centro-Oeste foi aquela onde a média, nessa questão, foi a maior ( 45,9 ), e a de menor média foi a região Sul ( 33,7 ).
Quanto à variabilidade das notas, o desvio padrão de todo o Brasil foi 33,1, inferior ao obtido na questão de número 1 ( 38,3 ).
O maior desvio nessa questão foi obtido na região Sudeste ( 33,5 ), enquanto o menor foi obtido na região Centro-Oeste ( 32,3 ).
A mediana do Brasil como um todo e das regiões Nordeste e Sudeste foi 40,0.
A maior mediana ocorreu na região Norte ( 47,5 ) e a menor na região Sul ( 35,0 ).
A nota máxima ( 100,0 ) e a nota mínima ( 0,0 ) foram obtidas em todas as regiões do Brasil.
O Gráfico 3.18 mostra a distribuição das notas na questão discursiva 2 do componente de Formação Geral para alunos do Bacharelado e da Licenciatura.
Observa -se que a maior frequência corresponde aos alunos que deixaram a questão em branco, tanto para aqueles de Bacharelado quanto para os de Licenciatura.
Aparece também um máximo local no intervalo ( 50 ; 60 ] para ambas as habilitações.
O coeficiente de assimetria da distribuição de notas dos concluintes de Licenciatura é negativo e bastante próximo de zero : – 0,01.
Para o Bacharelado a assimetria é positiva e mais acentuada, com coeficiente 0,59.
Destaca -se que por região os coeficientes de assimetria da Licenciatura são sempre próximos de zero, variando de – 0,27 a 0,15.
Ao contrário disso, os coeficientes de assimetria das distribuições de notas dos alunos provenientes do Bacharelado são positivos e revelam assimetrias positivas mais acentuadas, chegando a 1,14, na região Norte.
A assimetria do Bacharelado é negativa apenas para a região Centro-Oeste ( – 0,39 ). 3.3.1.4 Comentários sobre a correção das respostas à Questão Discursiva 2 Diferentemente da questão 1, a questão 2, cujo tema aborda políticas públicas para a erradicação do analfabetismo e para o aumento da empregabilidade, permitiu mais amplitude na elaboração das respostas dos estudantes, uma vez que solicita a apresentação de uma proposta.
À falta de objetividade técnica na elaboração da questão – sobretudo no seu comando – correspondeu uma miríade de possibilidades interpretativas.
Registre -se que não raras foram as respostas que se utilizaram de lugares comuns e exortações religiosas/humanitárias/cívicas, contudo, na grande maioria dos casos, essas exortações foram usadas como um encerramento do texto e não comprometeram a resposta.
Muito frequente também foi a confusão entre política pública e política partidária, bem como alguns poucos, confundiam programas educacionais com programas educativos nas TVs.
Assim, destacaram -se como propostas/programas mais recorrentes : • prosseguimento das políticas já existentes, mas com o aumento dos investimentos, normalmente sugerindo a ampliação da bolsa família, da Educação de Jovens e Adultos ( EJA ) e da EaD ; • valorização do magistério, construção de escolas e melhoria das atuais ; • parceria do governo com empresas para manutenção de salas de aula visando a alfabetização de seus funcionários, oferecimento de estágios e redução da carga horária em troca de isenção de impostos ; • parcerias com igrejas e ONGs para criação de espaços de alfabetização ; • escolas itinerantes e alfabetizadores em domicílio, principalmente para pessoas com dificuldade de locomoção, como os idosos e deficientes físicos ; • erradicação do trabalho infantil ; • vinculação da bolsa família não apenas à frequência, mas também e, principalmente, ao resultado obtido pelo aluno na escola ; • revisão das políticas atuais, sendo a mais frequente a extinção da aprovação automática ; • críticas consistentes em relação a modelos didáticos considerados inadequados e desestimulantes para a educação não só de adultos, mas de pessoas de todas as idades.
Destacam -se, ainda, outras sugestões apresentadas : • Creches nas escolas onde os pais estão sendo alfabetizados. • Diminuição de duas horas na jornada de trabalho em empresas para funcionários não alfabetizados, para que possam frequentar a escola. • Campanhas educativas vinculadas aos meios de comunicação.  Aumento do número de escolas noturnas.  Formação específica para professores alfabetizadores.
Quanto à relação entre o analfabetismo e a empregabilidade, deve -se sublinhar que nem todos os estudantes estabeleceram claramente o vínculo entre essas duas situações sociais.
Alguns falaram separadamente de uma e de outra.
Mas a maioria fez referência à necessidade de estudo para “ conseguir um bom emprego com um bom salário ”.
Alguns, em menor número, estabeleceram de forma bastante interessante a questão histórica para a situação do Nordeste ; e, também, a relação entre a pessoa analfabeta/com pouca instrução e a desempregada/empregada em serviços mais pesados e pior remunerados, que não teriam condições de educar seus filhos que, por sua vez, também não teriam melhores oportunidades no mercado de trabalho, identificando a formação de um círculo vicioso e a necessidade de sua interrupção.
Muitos afirmaram que o analfabetismo não é o único responsável pelo desemprego e sim a má distribuição de renda.
Grande também foi o número de alunos que criticaram o resultado da pesquisa, afirmando que a mesma é enganosa, uma vez que considera alfabetizada a pessoa que “ desenha seu nome ”, sendo comum associarem a atual política de alfabetização com ganhos eleitorais.
Foi comum a resposta incompleta, em que o estudante fez apenas a análise das desigualdades/crítica do quadro apresentado, ou só apresentou proposta.
Alguns se limitaram a responsabilizar o governo referindo -se de forma bastante genérica à questão da “ educação ” e “ profissional ”.
Entre os equívocos que mais se repetiram, destacou -se a simples análise dos dados apresentados na tabela que consta do enunciado.
Alguns textos, inclusive, estavam corretos, sem que, no entanto, fosse respondida a pergunta.
Da mesma forma, quando os estudantes partiram para a segunda parte da questão, fizeram referência a vários programas já existentes ou simplesmente disseram que é muito importante que existam projetos para a educação.
Outro equívoco recorrente foi a análise da educação no Brasil como um todo.
Essa análise, apesar de correta, não se referia especificamente à questão do analfabetismo, o que, por conseguinte, levava a sugestões que não eram direcionadas à erradicação do analfabetismo, e sim à melhora da educação no Brasil.
Sendo assim, foram apresentadas propostas como o aumento do número de faculdades, o reforço de alunos do Ensino Médio, a criação de escolas técnicas, etc.
Por outro lado, houve um grande número de redações bem escritas e precisas, no que se refere ao que foi exigido pela questão.
Foram análises equilibradas e sensatas, correlacionando a problemática do analfabetismo com o desemprego e a apresentação de sugestões bastante consistentes. 3.3.2 Componente de Conhecimento Específico Na parte da prova relativa às questões discursivas no componente de Conhecimento Específico ( Tabela 3.13 ), observa -se que a média foi mais baixa do que para as questões discursivas do componente de Formação Geral.
Enquanto no componente de Formação Geral a média para estudantes de Física de todo o Brasil foi 43,3, na parte de Conhecimento Específico a média foi 5,0.
A maior média deste componente foi obtida pelos estudantes da região Sudeste ( 6,7 ), e a menor, pelos da região Norte ( 2,2 ).
Quanto à variabilidade das notas, o desvio padrão de todo o Brasil foi 9,9.
O maior desvio padrão foi encontrado na região Sudeste ( 12,0 ), e o menor, na região Norte ( 5,1 ).
A maior nota máxima foi obtida na região Sudeste ( 68,3 ), enquanto a menor nota máxima foi encontrada na região Norte ( 33,3 ).
A nota mínima obtida por alunos de todas as regiões do Brasil, sem exceção, foi 0,0.
Além disso, a mediana do Brasil como um todo e em todas as regiões foi 0,0, o que indica que mais da metade dos alunos receberam nota zero nas questões discursivas de Conhecimento Específico de Física.
O Gráfico 3.19 representa a distribuição das notas nas questões discursivas no componente de Conhecimento Específico segundo a opção de Bacharelado/Licenciatura.
As modas destas distribuições ocorrem no intervalo, [ 0 ; 10 ].
O percentual de alunos de Licenciatura com notas no intervalo [ 0 ; 10 ] é maior do que o dos alunos de Bacharelado, enquanto esta proporção se inverte para todos os outros intervalos.
Nota -se que ambas as distribuições são decrescentes até o intervalo ( 60 ; 70 ], onde já ocorrem pouquíssimas notas.
Não se apresentam notas acima de 70,0.
Este comportamento das distribuições, assimetria positiva acentuada, é confirmado pelos coeficientes de assimetria positivos e maiores do que 1 para os dois cursos e em todas as regiões.
No Gráfico 3.19 os coeficientes de assimetria são 1,91 e 3,01, respectivamente, para Bacharelado e Licenciatura.
A análise de cada uma destas questões discursivas será feita a seguir. 3.3.2.1 Análise da Questão Discursiva 3 do Componente de Conhecimento Específico Na questão 3, cujos resultados aferidos encontram -se descritos na Tabela 3.14, a média dos estudantes de todo o Brasil foi 6,1.
A menor média nessa questão foi obtida pelos alunos da região Norte ( 4,3 ), enquanto a maior média foi obtida na região Sudeste ( 7,5 ).
Quanto à variabilidade das notas, o desvio padrão de todo o Brasil foi 14,5.
O maior desvio padrão foi obtido na região Sudeste ( 16,8 ), enquanto o menor foi obtido na região Norte ( 11,7 ).
A nota máxima, 100,0 pontos, foi alcançada por pelo menos um aluno da região Sudeste.
A região onde ocorreu a menor nota máxima foi a Nordeste ( 70,0 ).
A nota mínima ( 0,0 ) foi obtida em todas as regiões do Brasil, sem exceção.
As medianas de todas as regiões foram também zero, indicando que mais de metade dos estudantes de cada região receberam nota zero nesta questão.
O Gráfico 3.20 mostra a distribuição das notas na questão discursiva 3, do componente de Conhecimento Específico da área de Física para alunos do Bacharelado e da Licenciatura.
Ambas as distribuição têm moda nas questões em branco, com percentual um pouco maior para os alunos de Licenciatura ( compensado pelo percentual de nota zero entre os que responderam a questão ).
O coeficiente de assimetria é maior do que 2 para os dois cursos : 2,71 para o Bacharelado e 2,91 para a Licenciatura.
A assimetria é positiva e bastante acentuada para todas as distribuições por região, sendo o coeficiente de assimetria da distribuição das notas dos alunos de Licenciatura da região Norte, o maior ( 3,69 ). 3.3.2.2 Comentários sobre a correção das respostas à Questão Discursiva 3 O percentual de provas em branco, anuladas e desconsideradas foi bastante elevado.
Muitos estudantes deixaram a questão em branco, o que faz levantar a hipótese de que a maioria não assimilou os conhecimentos básicos de termodinâmica requeridos para a questão, ou não respondeu por desconsiderar o ENADE como instrumento válido de avaliação.
Grande parte dos concluintes teve, também, dificuldade para entender o que estava sendo solicitado, ou mesmo o que havia sido enunciado no texto da questão.
Os enunciados afirmam a impossibilidade de se construir uma máquina térmica ( Kelvin-Plank ) e um refrigerador ( Clausius ) com 100 \% de aproveitamento, isto é, sem perdas de energia ou necessidade de trabalho externo.
Outros tentaram exemplificar os enunciados sugerindo como exemplos dispositivos que violariam o afirmado por Kelvin-Plank e Clausius, utilizando muitas das vezes diagramas.
Tal tipo de resposta mostra que, apesar de algum contato com a termodinâmica, o estudante teve dificuldade de entender os enunciados de Kelvin-Plank.
Nesse caso, não houve percepção de que suas máquinas violariam as leis da Termodinâmica, o que resultou em baixa pontuação.
As respostas de pontuação média foram, na maioria das vezes, textos incompletos ou parcialmente corretos, onde os itens ( a ) e ( b ) foram respondidos, mas não o terceiro item.
O estudante entendeu o enunciado da questão, mas ainda apresentou alguma dificuldade em relação à afirmação de Clausius e, consequentemente, não respondeu ao terceiro item.
Os textos de pontuação mais elevada, muito poucos por sinal, responderam corretamente aos itens ( a ) e ( b ) e esboçaram alguma resposta para o item ( c ).
As leis da Termodinâmica são importantes para a Física e fazem parte da formação básica de um bom físico.
Contudo, os estudantes demonstraram muita dificuldade em entender aos enunciados e ao que foi solicitado na questão.
Mesmo aqueles com pouco tempo de exposição ao tema deveriam ser capazes de resolver a questão se tivessem entendido os enunciados de Kelvin-Plank e Clausius.
Isto pode significar que a maioria dos concluintes apresenta dificuldades de interpretação e compreensão de um texto, além de dificuldades na Física em si. 3.3.2.3 Análise da Questão Discursiva 4 do Componente de Conhecimento Específico A Tabela 3.15 contém as informações relativas à questão 4 do conjunto de questões do componente de Conhecimento Específico.
O desempenho dos estudantes de todo o Brasil nesta questão foi ainda mais baixo do que o desempenho na questão de número 3.
A média geral do Brasil foi 5,3, sendo a menor média registrada na região Norte ( 2,1 ), e a maior, na região Sul ( 7,2 ).
A nota máxima ( 100,0 ) foi atingida por pelo menos um aluno da região Sudeste.
A mediana e a nota mínima foram 0,0 em todas as regiões.
O Gráfico 3.21, representa a distribuição de notas na questão discursiva 4, no componente de Conhecimento Específico para alunos do Bacharelado e da Licenciatura.
A distribuição é bem semelhante à observada no Gráfico 3.20.
A maior frequência, tanto para estudantes de Bacharelado quanto para Licenciatura, ocorre na categoria “ em branco ”.
Observa -se que na categoria em branco e nos dois primeiros intervalos há uma proporção maior de alunos de Licenciatura, diferença compensada nos intervalos subsequentes.
Os coeficientes de assimetria são positivos e bastante altos : 2,25 para a distribuição de notas dos concluintes de Bacharelado e 3,44 para os de Licenciatura.
Este comportamento se repete para ambos os cursos em todas as regiões, sendo o menor coeficiente de assimetria igual a 1,98, da distribuição de notas de Bacharelado da região Sul, e o maior, 4,98, da Licenciatura da região Norte. 3.3.2.4 Comentários sobre a correção das respostas à Questão Discursiva 4 A questão 4 se referia a um dos pilares da Física Moderna : a dualidade onda-partícula, conceito que o estudante tem acesso ainda na formação básica de um curso de Física e, portanto, um tema de grau de dificuldade não elevado.
A questão estava suficientemente clara em seu enunciado.
No entanto, a quantidade de provas em branco, anuladas e desconsideradas foi bastante elevada.
Fundamentalmente, o padrão de resposta deve conter : discussão da dualidade da luz - comportamento de onda em fenômenos de difração e interferência e de partícula na interação com a matéria ( como no efeito fotoelétrico ou Compton ) ; discussão da dualidade da matéria ( no caso, elétron ) como postulado por de Broglie – partículas têm associado um comprimento de onda e podem difratar ou sofrer efeito de interferência ; finalmente, uma discussão dos experimentos apresentados no enunciado : tanto para a luz como para elétrons, ambos os comportamentos corpuscular e ondulatório ficam evidenciados, a distribuição de notas dos concluintes de Bacharelado e 3,44 para os de Licenciatura.
Este respectivamente, na interação com a tela de detecção e no padrão de interferência ( visível após acúmulo de suficiente estatística ).
Em relação às respostas em branco, é difícil avaliar se ocorreram devido ao desconhecimento do tema, falta de tempo ou falta de empenho durante o exame.
Entre os que responderam, a média, infelizmente, foi bastante baixa – e isto de fato é um pouco alarmante, dada a abrangência e relativa simplicidade dos conceitos abrangidos na questão.
Houve uma fração não desprezível de estudantes que sequer aludiu ao conceito da dualidade onda-partícula e, por isso, recebeu grau zero.
Por exemplo, houve vários que apenas voltaram a escrever o que já estava mencionado no enunciado.
Também houve casos de respostas simplesmente com conteúdo incorreto e/ou desconexo, em confronto com o enunciado.
Houve uma grande quantidade de notas baixas, porém não nulas ( entre 0,5 e 25,0 ), em que os estudantes fizeram algum tipo de alusão ao conceito solicitado.
Mas o fizeram de maneira muito imprecisa, inserindo erros, ou de forma bastante incompleta.
Pontuações um pouco mais altas ( entre 30,0 e 60,0 ) foram atribuídas àqueles que escreveram conceitos corretos, porém incompletos : mencionaram somente ao caráter dual da luz ou do elétron ( e não de ambos ), ou não aprofundaram o suficiente, ou não se referiram ao resultados em si.
Finalmente, graus mais altos ( \&gt ; 70,0 ) foram atribuídos a uma ( lamentavelmente ) pequena minoria que desenvolveu muito bem a questão, e veio a perder algum ponto apenas por deixar algumas lacunas.
De forma geral, a impressão é que, apesar de se tratar de uma questão com enunciado claro sobre um tema de formação essencial para um formando em Física, as respostas foram, salvo exceções, bastante fracas.
A maioria não se deu conta da sequência temporal das fotos que descreviam os experimentos, donde muitos deduziram que eram experimentos distintos. 3.3.2.5 Análise da Questão Discursiva 5 do Componente de Conhecimento Específico A Tabela 3.16 contém as informações relativas à questão 5 do conjunto do componente de Conhecimento Específico.
O desempenho dos estudantes nessa questão foi inferior ao das questões 3 e 4.
A nota média dos estudantes de todo o Brasil foi 3,6.
A maior média foi registrada na região Sudeste ( 5,9 ), enquanto a menor média foi registrada na região Norte ( 0,2 ).
Quanto à variabilidade das notas, o desvio padrão dos alunos do Brasil, como um todo, foi 12,4.
Enquanto o maior desvio foi encontrado na região Sudeste ( 15,6 ), o menor foi encontrado na região Norte ( 1,9 ).
O Gráfico 3.22 representa a distribuição das notas da questão discursiva 5 no componente de Conhecimento Específico para alunos do Bacharelado e da Licenciatura.
Mais uma vez, destaca -se o grande número de estudantes que deixaram a questão em branco, para duas modalidades de curso.
Como nas questões discursivas do Componente de Conhecimento Específico anteriores, esta tem assimetria positiva acentuada.
Para alunos do Bacharelado o coeficiente de assimetria é de 2,50 e para Licenciatura de 5,29.
Este tipo de assimetria também ocorre em todas as regiões, com coeficientes sempre maiores do que 2 para ambos os cursos.
A distribuição de notas da Licenciatura na região Norte tem o maior coeficiente de assimetria ( 11,59 ) e o menor é o das notas do Bacharelado dos alunos da região Sudeste ( 2,01 ). como um todo, foi 12,4.
Enquanto o maior desvio foi encontrado na região Sudeste ( 15,6 ), o menor foi encontrado na região Norte ( 1,9 ).
A nota máxima ( 100,0 ) foi alcançada na região Sudeste e a menor nota máxima foi registrada na região Norte ( 25,0 ).
A mediana e a nota mínima foram 0,0 para todas as regiões do Brasil, ou seja, pelo menos metade dos alunos tirou nota zero na questão em cada região.
O Gráfico 3.22 representa a distribuição das notas da questão discursiva 5 no componente de Conhecimento Específico para alunos do Bacharelado e da Licenciatura.
Mais uma vez, destaca -se o grande número de estudantes que deixaram a questão em branco, para duas modalidades de curso.
Como nas questões discursivas do Componente de Conhecimento Específico anteriores, esta tem assimetria positiva acentuada.
Para alunos do Bacharelado o coeficiente de assimetria é de 2,50 e para Licenciatura de 5,29.
Este tipo de assimetria também ocorre em todas as regiões, com coeficientes sempre maiores do que 2 para ambos os cursos.
A distribuição de notas da Licenciatura na região Norte tem o maior coeficiente de assimetria ( 11,59 ) e o menor é o das notas do Bacharelado dos alunos da região Sudeste ( 2,01 ). 3.3.2.6 Comentários sobre a correção das respostas à Questão Discursiva 5 A questão 5 versa sobre um sistema massa-mola.
Os conceitos físicos e descrições matemáticas para o sistema massa-mola são extensamente estudados ao longo do curso de Física, visto que representam o protótipo de movimentos harmônicos simples.
Assim sendo, o tema abrangido nesta questão é de caráter de formação geral de um físico.
Entretanto, a quantidade de provas em branco, anuladas e não consideradas foi a mais elevada dentre as três discursivas do Componente de Conhecimento Específico.
A questão apresentava um gráfico de resposta de frequência H(ω) versus frequência angular ω.
Era também fornecida a expressão analítica para a função H(ω) : Constava de 4 subitens, cada um deles valendo 2,5 pontos.
Os itens ( a ) e ( b ) eram mais imediatos, pois requeriam apenas uma leitura correta do gráfico : o item ( a ) pedia a frequência de ressonância ωr e bastava ver, no gráfico, para qual valor de ωr o valor de H(ω) era máximo ; o item ( b ) pedia o valor da constante elástica da mola, k, que também podia ser obtida diretamente do gráfico já que H(ω = 0 ) = 1/k.
Os itens ( c ) e ( d ) eram de um grau de dificuldade um pouco mais elevado : pediam, respectivamente, os valores de c e m e, para isso, era necessário dar -se conta que, sob aproximação, ωr ≈√k/m e assim Hmax = H(ωr) = 1/(cωr) ( para o item (c)) e m≈kωr2 ( para o item ( d ) ).
A maior surpresa com relação às respostas dos alunos foi o excessivo percentual de respostas em branco.
Algumas tentativas de diagnóstico para este alto índice podem ser : cansaço e/ou falta de tempo por ser última questão da prova ; dificuldade com o caráter gráfico e numérico da questão.
Daqueles que tentaram desenvolver a questão, a média ficou um pouco acima dos 20,0 pontos, refletindo o fato de que a maioria dos alunos conseguiu ao menos fazer o item ( a ) ou o item ( b ) ; entre estes, houve uma divisão entre os que tentaram resolver os outros itens ( e erravam ) ou que simplesmente deixavam em branco o restante.
Pontuações médias ( entre 30,5 e 60,0 ), salvo exceções, foram, em geral, devidas aos estudantes que conseguiram desenvolver os itens ( a ) e ( b ), mas erraram ou deixaram em branco os itens ( c ) e ( d ).
Pontuações altas ( acima de 70,0 ) foram obtidas tipicamente por aqueles que tentaram e acertaram ao menos três dos quatro itens.
Houve vários alunos que acertaram ( numericamente ) todas as respostas, porém com omissões nas justificativas.
Ou então, que elaboraram todos os itens corretamente, mas faltou justificativa da aproximação nos itens ( c ) e ( d ) : O altíssimo percentual de respostas em branco de certa forma prejudica a análise global/estatística do resultado da questão.
Daqueles que responderam a questão, houve bastante acerto para os itens ( a ) e/ou ( b ), que eram de grau de dificuldade baixo, e apenas uma minoria desenvolveu de forma parcial ou totalmente correta os itens ( c ) e ( d ) ( que já não eram tão fáceis ). 3.3.3 Considerações Finais da Banca de correção das questões discursivas do Componente de Conhecimento Específico Em conclusão, a correção das questões discursivas do ENADE 2011 de Física evidenciou grande despreparo dos estudantes que nele participaram.
Dada a ressalva que talvez uma parcela dos concluintes mais preparados pode haver escolhido não participar da avaliação ao entregar a prova em branco, o que é apenas uma especulação sem nenhuma evidência no momento, os resultados em geral mostram grandes dificuldades para lidar com questões conceituais básicas, e dificuldades maiores ainda para atacar problemas que envolvam manipulações matemáticas nem tão complexas assim.
Em especial, sobre a questão de mecânica quântica, a falta de conhecimento sobre a dualidade onda-partícula é alarmante.
Estudantes em final de curso não deveriam cometer os erros, enganos e mistificações observadas durante as correções.
Este exame é uma oportunidade, a ser aperfeiçoada nas edições posteriores, de \textbf{traçar} um painel acurado da formação básica atual do jovem professor/pesquisador em Física no Brasil.
Pelo resultado deste ano, esta formação, em média, leva a resultados pouco encorajadores, o que deve ser ajustado de alguma maneira.
Apesar de tudo, excelentes alunos participaram ativamente do ENADE 2011 e mostraram que existe um número razoável de concluintes que consegue ter uma performance muito boa, e que pode ser capaz de “ carregar a tocha ” para as futuras gerações.

% matched lemmas: traçar
\end{textsample}
